
\section{Preliminaries}\label{sec:aggregation}
We fix the semantics of finite-variable formulas, relational terms, and
circuits, with examples that will be used throughout the paper. All
structures have nonempty domains. Our lower bounds concern finite
structures, whereas the explicit translations hold on all nonempty
structures.

\subsection{Finite-variable logic and binary queries}\label{sec:prelim-logic}
\paragraph{Signatures and formulas.}
A \emph{signature} $\sigma$ is a finite set of binary relation symbols;
there are no function or individual constant symbols. A $\sigma$-structure
$\mathcal M$ has domain $D\coloneqq\operatorname{dom}(\mathcal M)\ne\varnothing$
and interprets each $R\in\sigma$ as a relation $R^{\mathcal M}\subseteq D^2$.
Equality is available. For $\ell\ge1$, the \emph{$\ell$-variable fragment}
$\mathrm{FO}^{\ell}[\sigma]$ consists of first-order formulas using only
$\ell$ variable names, with arbitrary reuse and rebinding of those names.
Formally, with $u,v$ ranging over a fixed set of $\ell$ names, its syntax is
\[
 \varphi::=R(u,v)\mid u=v\mid\neg\varphi
 \mid(\varphi\land\varphi)\mid(\varphi\lor\varphi)
 \mid\exists u\,\varphi\mid\forall u\,\varphi.
\]
We use the names $x,y$ for \FOtwo\ and $x,y,z$ for \FOthree.
The restriction is on the \emph{number of variable names}, not on the
number or nesting depth of quantifiers. For example,
\[
 \varphi(x,y)\coloneqq P(x,y)\land\exists y\,Q(x,y)
\]
is an \FOtwo\ formula: the quantifier rebinds $y$ only in its own scope,
leaving the occurrence of $y$ in $P(x,y)$ free. Write
$\operatorname{FV}(\varphi)$ for the set of free variables.

\paragraph{Assignment relations.}
A formula with $\operatorname{FV}(\varphi)\subseteq\{x,y\}$ defines the
\emph{binary query}
\[
 \sem\varphi{\mathcal M}
 \coloneqq\{(a,b)\in D^2:\mathcal M\models
       \varphi[x\mapsto a,y\mapsto b]\}.
\]
A variable absent from $\operatorname{FV}(\varphi)$ contributes a
\emph{dummy coordinate}: the query is independent of its assigned value.
Thus a sentence $\varphi$ has the binary assignment value
\[
 \sem\varphi{\mathcal M}
 =\begin{cases}
     D^2,&\mathcal M\models\varphi,\\
     \varnothing,&\mathcal M\not\models\varphi.
   \end{cases}
\]
This convention lets formulas, relational terms, and circuits have the
same kind of output. An \FOthree\ subformula may nevertheless have
\emph{three} free variables: for example,
$\theta(x,y,z)\coloneqq R(x,z)\land S(z,y)$ has the ternary assignment
relation $\{(a,b,c)\in D^3:R(a,c)\land S(c,b)\}$. Quantifying $z$
returns a binary query. Accordingly, direct \FOthree\ evaluation can
store three-variable truth tables in $O(|\varphi|n^3)$ time on an
$n$-element domain, whereas every intermediate value of a \CoR\ term
or circuit is binary. Section~\ref{sec:matrix} revisits this arity
contrast for matrix query languages.

\subsection{The calculus of relations and relational circuits}\label{sec:prelim-cor}
\paragraph{Relational terms.}
The \emph{calculus of relations} (\CoR) has terms generated by
\[
 T::=R\mid\zero\mid\one\mid\id
   \mid(T\cup T)\mid(T\cap T)\mid\overline T
   \mid T^\smile\mid(T;T), \qquad R\in\sigma.
\]
Every term denotes a binary relation on $D$. Its constants are
$\zero\coloneqq\varnothing$, $\one\coloneqq D^2$, and
$\id\coloneqq\{(a,a):a\in D\}$. Union and intersection are set operations;
for relations $X,Y\subseteq D^2$ the other operations are
\[
  \overline X \coloneqq D^2\setminus X,\quad
  X^\smile \coloneqq \set{(b,a):(a,b)\in X},\quad
  (X;Y) \coloneqq 
     \set{ (a, b) : \exists c\in D. (X(a,c)\land Y(c,b)) }.
\]
Here $X(a,b)$ abbreviates $(a,b)\in X$; complement is unrestricted,
including around compositions. A \emph{composition} $X;Y$ is Boolean
matrix multiplication, with $c$ its existential \emph{witness}.

Two useful special cases are the \emph{row and column projections}:
\begin{equation*}\tag{Proj}\label{eq:prelim-projections}
 \begin{aligned}
  (X;\one)(a,b)&\iff\exists c\in D\ X(a,c),\\
  (\one;X)(a,b)&\iff\exists c\in D\ X(c,b).
 \end{aligned}
\end{equation*}
The first is independent of $b$, and the second of $a$. More generally,
for binary predicates $A,B,C$ with the indicated variable pairs,
\begin{equation*}\tag{QE}\label{eq:prelim-elimination}
 \exists z\,\bigl(A(x,y)\land B(x,z)\land C(z,y)\bigr)
 \quad\equiv\quad \bigl(A\cap(B;C)\bigr)(x,y).
\end{equation*}
For example, the \FOtwo\ query
$P(x,y)\land\exists y\,Q(x,y)$ above is $P\cap(Q;\one)$ in \CoR.
The position of complement matters: $\overline{R;\one}$ tests that
$x$ has no $R$-successor, whereas $\overline R;\one$ tests that
$x$ has some non-$R$ pair. These queries are not equivalent.

\paragraph{A linear translation from \FOtwo.}\label{sec:FOtwo-translation}
Under the assignment convention above, the standard translation
$\tau$ from \FOtwo\ into \CoR\ is given on atoms and quantifiers by
\[
\begin{array}{c|c@{\qquad}c|c}
 \varphi&\tau(\varphi)&\varphi&\tau(\varphi)\\\hline
 R(x,y)&R&R(y,x)&R^\smile\\
 R(x,x)&(R\cap\id);\one&R(y,y)&\one;(R\cap\id)\\
 x=y\ \text{or}\ y=x&\id&x=x\ \text{or}\ y=y&\one\\
 \exists x\,\psi&\one;\tau(\psi)&
 \exists y\,\psi&\tau(\psi);\one
\end{array}
\]
Conjunction, disjunction, and negation become intersection, union, and
complement; universal quantifiers use
$\forall u\,\psi\equiv\neg\exists u\,\neg\psi$. Structural induction,
using~\eqref{eq:prelim-projections}, shows that
\[
 \sem{\tau(\varphi)}{\mathcal M}
   =\sem\varphi{\mathcal M},\qquad
 |\tau(\varphi)|=O(|\varphi|).
\]
The translation also preserves the number of relation-atom
occurrences, including reflexive atoms. See
\citep[Remark~3.22]{nakamura2022}. Section~\ref{sec:FOtwo} bounds the fold cost of this translation
by its number of quantifiers.

\paragraph{Terms versus circuits.}
A \emph{term} is a syntax tree, so repeated occurrences of a subterm
are distinct nodes. A \emph{relational circuit} is a finite directed
acyclic graph (DAG): each gate is an input, a constant, or one of the
same \CoR\ operations, with ordered input ports, and a designated
output gate. A gate may feed multiple later gates, including both
ports of one composition. Its value $\sem g{\mathcal M}$ is computed
once from its predecessors; $\sem C{\mathcal M}$ is the value of the
output gate.

For example, the term $(R;R);(R;R)$ has seven nodes and three
composition occurrences. The circuit
\[
 g_0\coloneqq R,\qquad g_1\coloneqq g_0;g_0,\qquad g_2\coloneqq g_1;g_1
\]
has just three distinct gates, two of which are compositions
(Figure~\ref{fig:term-vs-circuit}).

\begin{figure}[H]
\centering
\begin{tikzpicture}[
  comp/.style={draw,circle,fill=blue!7,minimum size=7mm,inner sep=0pt},
  inp/.style={draw,rounded corners=2pt,fill=gray!8,minimum width=8mm,
              minimum height=6mm,inner sep=1pt},
  every node/.style={font=\small},>=Latex]
 \begin{scope}[xshift=-2.25cm]
   \node[font=\small\bfseries] at (0,0.65) {Term (tree)};
   \node[comp] (t0) at (0,0) {$;$};
   \node[comp] (t1) at (-0.95,-0.9) {$;$};
   \node[comp] (t2) at (0.95,-0.9) {$;$};
   \node[inp] (t3) at (-1.45,-1.8) {$R$};
   \node[inp] (t4) at (-0.45,-1.8) {$R$};
   \node[inp] (t5) at (0.45,-1.8) {$R$};
   \node[inp] (t6) at (1.45,-1.8) {$R$};
   \draw (t0)--(t1) (t0)--(t2);
   \draw (t1)--(t3) (t1)--(t4) (t2)--(t5) (t2)--(t6);
   \node[font=\footnotesize] at (0,-2.30) {3 composition nodes};
 \end{scope}
 \begin{scope}[xshift=2.25cm]
   \node[font=\small\bfseries] at (0,0.65) {Circuit (DAG)};
   \node[comp] (c2) at (0,0) {$;$};
   \node[comp] (c1) at (0,-0.9) {$;$};
   \node[inp] (c0) at (0,-1.8) {$R$};
   \draw[-{Latex}] (c0.north west) to[out=130,in=225] (c1.south west);
   \draw[-{Latex}] (c0.north east) to[out=50,in=315] (c1.south east);
   \draw[-{Latex}] (c1.north west) to[out=130,in=225] (c2.south west);
   \draw[-{Latex}] (c1.north east) to[out=50,in=315] (c2.south east);
   \node[font=\footnotesize] at (0,-2.30) {2 composition gates};
 \end{scope}
\end{tikzpicture}
\caption{The term $(R;R);(R;R)$ repeats the same subterm, whereas
 the circuit reuses one gate as both inputs to the next composition.
 Distinct gates, not wires or uses, are counted.}
\label{fig:term-vs-circuit}
\end{figure}

More generally, $g_d$ defined by
$g_{i+1}\coloneqq g_i;g_i$ computes the pairs joined by an $R$-walk of length
$2^d$. Fully unfolding this computation gives a term with $2^d-1$
compositions, while the circuit has only $d$ composition gates.
For another example, a circuit can compute $h\coloneqq R;S$ once and use it
in both branches of $(P\cap h)\cup(Q\cap h^\smile)$, still with only
one composition gate. Thus sharing can save exponentially many
compositions; our lower bounds hold \emph{despite} allowing it.

\paragraph{Equations.}
A \emph{\CoR\ equation} is a statement $T=S$, satisfied by a
structure precisely when
\[
 \mathcal M\models T=S
  \quad\iff\quad
 \sem T{\mathcal M}=\sem S{\mathcal M}.
\]
For example, $\id\cup R=\one$ expresses the \FOtwo\ sentence
$\forall x\,\forall y\,(x=y\lor R(x,y))$. A circuit computing an
equation may have two designated output gates (for $T$ and $S$),
with shared intermediate gates counted only once.

\subsection{Resource measures and computational conventions}\label{sec:prelim-resources}
\paragraph{Syntax size and gate counts.}
Formula size $|\varphi|$ and term size $|T|$ count nodes in their
\emph{syntax trees}; all connectives have arity at most two. Negation
and quantification are unary, and displayed iterated connectives are
abbreviations to be expanded. Each relation-atom occurrence counts as
one node even for a growing signature (so $|G_k|=8k$); the length of
the written relation name is not charged. This is different from
measuring the encoded bit length of a formula. Write
$\relatoms\varphi$ for the number of relation-atom occurrences
(including $R(x,x)$, but excluding equality atoms), and
$\#_{\exists\forall}(\varphi)$ for its quantifier occurrences.

For a term $T$, write $\relatoms T$ for occurrences of input relation
symbols and $\cost T$ for occurrences of composition. For a circuit
$C$, $|C|$ counts its distinct gates, including input and constant
gates, and $\gcost C$ counts distinct composition gates. All other
operations are uncharged by $\gcost C$, though they do count toward
$|C|$. A shared gate counts once, irrespective of its fan-out. We will
also use the elementary tree inequalities
\[
 |T|\ge2\cost T+1,
 \qquad |T|\ge2\relatoms T-1.
\]

\paragraph{Equivalence and allowed inputs.}
For a formula $\varphi$ and a relational term or circuit $C$ over
the same signature, write
\[
 C\eqfin\varphi
 \quad\iff\quad
 \sem C{\mathcal M}=\sem\varphi{\mathcal M}
 \quad\text{for every finite $\sigma$-structure $\mathcal M$.}
\]
The same notation applies to terms and to comparisons of formulas.
Unsubscripted equivalences in our explicit translations mean
agreement over all nonempty structures. A \emph{size-specific} circuit
$C_n$ may depend on $n=|D|$, but not on the input interpretation;
it must compute the intended query on \emph{every} $n$-element
input. A relational circuit has only the input symbols and the three
constants $\zero,\one,\id$; no built-in order, named vertices,
extra vertex predicates, or arbitrary vertex-dependent constants are
supplied for free. The more permissive typed matrix circuit model of
Section~\ref{sec:typed} is defined separately.

\paragraph{Randomized targets.}
A randomized circuit is an input-independent distribution
$\mathcal D_n$ over deterministic circuits, possibly of different
sizes, for the given domain size $n$. Pointwise error at most
$\varepsilon$ for a query $q$ means that, for every $n$-element input
$\mathcal M$ and every $(a,b)\in D^2$,
\[
 \Pr_{C\sim\mathcal D_n}\!\left[
   \sem C{\mathcal M}(a,b)=\sem q{\mathcal M}(a,b)
 \right]\ge1-\varepsilon.
\]
The entire output relation need not be correct simultaneously.
Randomized resource bounds concern expectations over $\mathcal D_n$;
the circuit distribution is fixed independently of the input.

\paragraph{Common notation.}
We use $[k]\coloneqq\{1,\ldots,k\}$ and $2^{[k]}$ for its power set;
$A^c\coloneqq[k]\setminus A$ when $A\subseteq[k]$. Empty intersections and
unions of relations mean $\one$ and $\zero$, respectively; the
analogous Boolean operations mean true and false. Truth values are
$\bool\coloneqq\{0,1\}$, and $[\mathcal P]$ is the indicator of a statement
$\mathcal P$. The special palette, hub, and fold-prediction notation
are introduced when needed in Section~\ref{sec:preservation}.
